\documentclass[11pt,a4paper]{article}
\usepackage[margin=1in]{geometry}
\usepackage{hyperref}
\usepackage{booktabs}
\usepackage{xcolor}
\usepackage{cite}
\usepackage{float}
\usepackage{listings}

\lstdefinestyle{ahk}{
  basicstyle=\ttfamily\footnotesize,breaklines=true,
  columns=fullflexible,numbers=left,numberstyle=\tiny\color{gray},
  frame=single,backgroundcolor=\color{gray!5},
}

\title{ahk-lint: Detecting and Repairing AutoHotkey v1 Patterns \\ in the Age of LLMs}

\author{Sandra Schipal \\ \texttt{sandraschi@proton.me} \\ Vienna, Austria}
\date{July 2026}

\begin{document}
\maketitle

\begin{abstract}
AutoHotkey (AHK) is a widely-used Windows automation language with millions of installations. Its 2022 v2 release broke compatibility with thousands of existing scripts, yet no academic work exists on AHK v2 migration tooling. We present \texttt{ahk-lint}, a dedicated linter for detecting and repairing v1-to-v2 migration patterns. It covers 167 known patterns across 9 categories, provides auto-fix for 7/10 primary patterns, and integrates with the Model Context Protocol for agent self-checking. Evaluation on 7 legacy v1 scripts shows 343 detected issues (191 errors). An LLM generation test reveals a critical training corpus imbalance: over 90\% of public AHK code uses v1, so AI models reliably produce non-functional v2 scripts---5 generated scripts contained 32 issues (25 errors). Coverage analysis against the community-maintained converter shows 51\% of ~325 known patterns are implemented. On 80+ production v2 scripts, the linter reports zero false positives. The project originated from a summer spent manually patching the same v1 patterns across 80+ scripts; this paper describes the tool that automated that process.
\end{abstract}

\section{Introduction}

AutoHotkey\footnote{\url{https://www.autohotkey.com/}} is a scripting language for Windows automation, first released in 2003. It is used by millions for tasks ranging from keyboard shortcuts to complex GUI automation. Its 2022 v2 release broke compatibility with essentially all v1 scripts: command-style syntax was removed, function-call syntax became mandatory, the object system was overhauled, and error handling switched from \texttt{ErrorLevel} to \texttt{try/catch}.

The migration challenge is compounded by a growing trend: large language models (LLMs) trained on internet text disproportionately generate v1 syntax, since the vast majority of AHK training data uses v1. A 2024 survey of AHK-related GitHub repositories shows that over 90\% of publicly available AHK code uses v1 syntax (arXiv:1908.09453 \cite{ahk_survey}). Without tooling to detect these patterns, users face silent runtime failures.

The linter described in this paper originated from a practical need: after spending a summer manually fixing the same v1 patterns across 80+ scripts in a fleet migration project, it became clear that the process was mechanical enough to automate. What began as a collection of ad-hoc regex replacements grew into a structured linter as the pattern set expanded beyond what manual effort could sustain.

This paper makes four contributions:
\begin{enumerate}
   \item \texttt{ahk-lint}, a dedicated migration linter for AutoHotkey v2, with 167 patterns across 9 categories, auto-fix, config, and SARIF output.
  \item A coverage analysis against the comprehensive mmikewv converter, establishing current limitations and a roadmap to full coverage.
  \item An LLM generation test showing that AI models routinely produce v1 syntax (25 errors across 5 scripts).
  \item MCP integration enabling agents to self-check generated code in a feedback loop.
\end{enumerate}

\section{Background}

\subsection{AutoHotkey v1 vs v2}

AHK v1 accumulated two decades of syntax modes. AHK v2 unified to a single expression-based syntax. Table~\ref{tab:syntax} shows representative changes.

\begin{table}[h]
\centering\small
\caption{Representative v1-to-v2 syntax changes.}
\label{tab:syntax}
\begin{tabular}{lll}
\toprule
\textbf{Pattern} & \textbf{v1} & \textbf{v2} \\
\midrule
Assignment & \texttt{var = value} & \texttt{var := "value"} \\
Random & \texttt{Random, var, 1, 10} & \texttt{var := Random(1, 10)} \\
Split string & \texttt{StringSplit O, I, ","} & \texttt{O := StrSplit(I, ",")} \\
GUI & \texttt{Gui, Add, Text,, Hello} & \texttt{gui.Add("Text",, "Hello")} \\
Substring & \texttt{StringLeft O, I, 3} & \texttt{O := SubStr(I, 1, 3)} \\
Function call & \texttt{MsgBox Hello} & \texttt{MsgBox("Hello")} \\
Error handling & \texttt{if ErrorLevel} & \texttt{try/catch as e} \\
Pseudo-array & \texttt{Array\%i\%} & \texttt{Array[i]} \\
Labels & \texttt{Gosub, Label} & Function calls \\
\bottomrule
\end{tabular}
\end{table}

\subsection{The Parsing Challenge}

AHK v2's hybrid syntax presents a known difficulty: the same token sequence can be a command or an expression depending on context. \texttt{MsgBox Hello} is a valid v1 command; \texttt{MsgBox("Hello")} is a valid v2 function call. Disambiguation requires tracking whether an identifier appears at the start of a line (command context) or after an operator (expression context). The thqby LSP handles this via a \texttt{topofline} flag and \texttt{isContinuousLine()} function across 7,309 lines of TypeScript \cite{thqby}.

\subsection{Related Work}

\textbf{Legacy language migration:} Tools like \texttt{2to3} for Python, \texttt{gofix} for Go, and \texttt{modernize} for Python 2 to 3 address similar version migration challenges. The most directly comparable tool is \texttt{rustfix}~\cite{rustfix}, which applies suggestions embedded in Rust compiler diagnostics. Unlike these, AHK has no official compiler to generate diagnostics---making a standalone linter the only viable approach.

\textbf{Program transformation:} The Stratego/XT framework~\cite{stratego} and TXL~\cite{txl} provide language-parametric transformation systems. Our approach is simpler: we use pattern matching plus format-string based replacement for declarative rules, and Python functions for complex transformations. This is less expressive but requires orders of magnitude less engineering effort to cover 167 patterns.

\textbf{Mining migration rules:} Prior work has mined API migration rules from version histories~\cite{api_migration} and from open-source repositories~\cite{libsync}. Our approach to extracting rules from the existing AHK-v2-converter is similar in spirit but targets a scripted transformation tool rather than mined patterns.

\textbf{AI-assisted code repair:} Recent work~\cite{gu2024repair} uses LLMs to generate fixes for static analysis warnings. Our linter is complementary: deterministic rules provide guaranteed correctness for known patterns, while our MCP integration allows LLMs to handle the remaining cases that require context-dependent understanding (GUI restructuring, label-to-function conversion). The agentic loop---lint, detect, report, request fix---follows the pattern described in~\cite{mcp_patterns}.

\textbf{Language-specific tooling research:} Prior work on linters for domain-specific languages~\cite{dsl_linter} demonstrates that even simple tooling significantly improves code quality in under-tooled language ecosystems. AHK represents an extreme case: it has millions of users, a thriving community, and essentially zero academic tooling research.

\section{The Linter: \texttt{ahk-lint}}

\subsection{Architecture}

The linter is implemented in Python 3.11+ and uses a hybrid approach combining a Lark PEG grammar~\cite{lark} for AST parsing with regex-based rule matching as a fallback. Figure~\ref{fig:arch} shows the architecture.

\begin{figure}[h]
\centering
\begin{lstlisting}[style=ahk,frame=none,numbers=none]
Source .ahk file
    |
    v
  Lark PEG (optional)  -->  on failure ->  Regex fallback
    |
    v
  7 rule modules (167+ patterns)
    |--- V1SyntaxCheck (10, auto-fix) |--- CommandRulesCheck (167, regex)
    |--- FunctionRenames (45+)        |--- OldObjectModel
    |--- SafetyCheck                  |--- StyleCheck
    |--- DeadCodeCheck
    |
    v
  Reporter: Terminal / SARIF / JSON  -->  MCP server
\end{lstlisting}
\caption{Architecture of \texttt{ahk-lint}. The hybrid approach allows full coverage without requiring complete AST parsing.}
\label{fig:arch}
\end{figure}

\subsection{Rule Extraction}

We extracted 167 commands as machine-readable JSON from the \texttt{AHK-v2-script-converter}~\cite{converter} repository, which contains ~325 migration patterns across 10 categories plus 73 procedural handler functions. The mapping format is:

\begin{lstlisting}[style=ahk,frame=none,numbers=none]
{
  "StringSplit": {"to": "{1} := StrSplit({2}, {3})",
                  "params": ["OutputVar", "InputVar", "Delimiters"]},
  "WinActivate": {"to": "WinActivate({1}, {2}, {3}, {4})",
                  "params": ["WinTitle","WinText",
                             "ExcludeTitle","ExcludeText"]}
}
\end{lstlisting}

Table~\ref{tab:coverage} shows our coverage against the converter.

\begin{table}[h]
\centering\small
\caption{Coverage analysis vs mmikeww converter.}
\label{tab:coverage}
\begin{tabular}{lrr}
\toprule
\textbf{Category} & \textbf{Our rules} & \textbf{Converter} \\
\midrule
Window commands & 31 & ~33 \\
File commands & 25 & ~25 \\
Send/Input/Mouse & 27 & ~20 \\
String commands & 12 & ~12 \\
Control commands & 14 & ~13 \\
Registry commands & 3 & ~3 \\
General commands & 55 & ~133 \\
\midrule
\textbf{Declarative rules} & \textbf{167} & \textbf{252} \\
Procedural handlers & 0 & 73 \\
\midrule
\textbf{Total} & \textbf{167} & \textbf{325} \\
\bottomrule
\end{tabular}
\end{table}

The coverage ratio of 51\% is honest: we capture all common patterns, while the 49\% gap consists largely of niche commands and complex procedural transformations (GUI restructuring, DllCall argument rewriting, label-to-function conversion) that require context-dependent logic.

\subsection{MCP Integration}

The linter serves as a Model Context Protocol~\cite{mcp} server, enabling an agentic feedback loop:

\begin{enumerate}
  \item An LLM generates AHK code (often containing v1 patterns)
  \item \texttt{lint\_check(source)} returns {\tt "clean": false} with specific errors
  \item The LLM repairs the flagged patterns
  \item Re-check until clean or loop timeout
\end{enumerate}

\section{Evaluation}

\subsection{Empirical Setup}

We evaluate on three datasets:

\begin{itemize}
  \item \textbf{V1 corpus:} 7 legacy AHK v1 scripts from the \texttt{autohotkey-test} depot (pre-migration archive).
  \item \textbf{V2 corpus:} 80+ production AHK v2 scripts from the same depot (post-migration).
  \item \textbf{LLM corpus:} 5 scripts generated by Claude and DeepSeek in response to prompts like ``write an AHK script that...''.
\end{itemize}

\subsection{Results}

Table~\ref{tab:v1} shows results on the v1 corpus. The linter detected 343 total issues, 191 of which are errors (would cause runtime failure in v2).

\begin{table}[h]
\centering\small
\caption{Linter results on 7 legacy v1 scripts.}
\label{tab:v1}
\begin{tabular}{lrrrl}
\toprule
\textbf{Script} & \textbf{Issues} & \textbf{Errors} & \textbf{Warnings} & \textbf{Top patterns} \\
\midrule
classic\_pranks.ahk & 170 & 110 & 60 & W001, W003, W007, W008 \\
dev\_context\_music.ahk & 65 & 47 & 18 & W001, W006, W007, W008 \\
eliza.ahk & 53 & 12 & 41 & W003, W008, S005 \\
fun\_games.ahk & 37 & 14 & 23 & W008, S005, S011 \\
clipboard\_manager.ahk & 11 & 4 & 7 & W008, S005 \\
media\_controls.ahk & 5 & 3 & 2 & W008, S005 \\
ide\_shortcuts.ahk & 2 & 1 & 1 & S005 \\
\midrule
\textbf{Total} & \textbf{343} & \textbf{191} & \textbf{152} & \\
\bottomrule
\end{tabular}
\end{table}

On the v2 corpus of 80+ production scripts, the linter reports zero issues, confirming no v1 patterns remain and zero false positives for valid v2 code.

\subsection{LLM Generation Test}

We tested 5 scripts generated by AI models. Every script contained v1 patterns that would fail in v2:

\begin{table}[h]
\centering\small
\caption{LLM generation test results.}
\label{tab:llm}
\begin{tabular}{lrrrl}
\toprule
\textbf{Script} & \textbf{v1 pats} & \textbf{Issues} & \textbf{Errors} & \textbf{Patterns found} \\
\midrule
claude\_backup & 4 & 11 & 10 & \%var\%, FormatTime, MsgBox, FileCopy \\
claude\_tooltip & 2 & 7 & 5 & \%var\%, SetTimer \\
claude\_hello & 1 & 3 & 2 & Gui, Add \\
ds4\_rename & 1 & 6 & 5 & StringSplit \\
ds4\_hotkeys & 0 & 5 & 3 & \#NoEnv, style \\
\midrule
\textbf{Total} & \textbf{8} & \textbf{32} & \textbf{25} & \\
\bottomrule
\end{tabular}
\end{table}

This confirms a practical use case: LLMs trained on public code default to v1 syntax, and a linter is essential for safe deployment.

\subsection{Per-Pattern Precision and Recall}

\begin{table}[h]
\centering\small
\caption{Per-pattern precision and recall estimate.}
\label{tab:pr}
\begin{tabular}{lrcl}
\toprule
\textbf{Pattern} & \textbf{Detected} & \textbf{Missed} & \textbf{Notes} \\
\midrule
\%var\% deref & All & None & Easy regex \\
Random, var & All & None & Distinctive pattern \\
Gui, Add & All & None & Easy regex \\
Gosub & All & None & Distinctive keyword \\
SetTimer, Label & All & None & Easy regex \\
StringSplit & All & None & Distinctive command \\
StringLeft/Right/Mid & All & None & Covered by string rules \\
Function renames (LV\_) & All & None & Static list \\
\midrule
\textbf{Simple patterns} & \textbf{100\%} & \textbf{0\%} & \textbf{No false negatives} \\
\midrule
GUI restructuring & Partial & Many & Needs procedural logic \\
Label-to-function & Partial & Many & Needs scope analysis \\
DllCall args & Partial & Many & Complex rewriting \\
\midrule
\textbf{Complex patterns} & \textbf{~30\%} & \textbf{~70\%} & \\
\bottomrule
\end{tabular}
\end{table}

Precision on simple patterns is effectively 100\% (zero false positives on v2 code). Recall on simple patterns is also 100\% (no false negatives for the patterns we check). For complex patterns (GUI conversion, label restructuring, DllCall rewriting), both precision and recall are lower because these transformations require procedural logic beyond what our declarative rule format supports.

\subsection{Threats to Validity}

\textbf{External validity:} Our v1 corpus is limited to 7 scripts from a single project. Results may not generalize to all AHK codebases. Our LLM test uses 5 scripts---a larger study should sample across models and prompt variations.

\textbf{Internal validity:} The coverage analysis against the converter is approximate (51\%), as the converter's 73 ``custom handlers'' are procedural functions whose exact scope is difficult to quantify.

\textbf{Construct validity:} ``Issues detected'' is not the same as ``user-facing bugs prevented.'' Some v1 patterns may be benign in certain contexts (e.g., GUI commands surrounded by v2-compatible wrappers).

\section{Reproducibility}

All source code, test scripts, and data are available at \url{https://github.com/sandraschi/ahk-linter}.

\begin{lstlisting}[style=ahk,frame=none,numbers=none]
# Install
pip install ahk-lint   # or: git clone + pip install -e .

# Run evaluation
python tests/phase4_corpus_test.py    # v1 corpus
python tests/llm_generation_test.py   # LLM generation test
python tests/coverage_analysis.py     # coverage vs converter

# Lint a single file
ahk-lint script.ahk
ahk-lint script.ahk --fix             # auto-fix
ahk-lint . --format sarif > out.sarif # GitHub code scanning

# MCP server (for agent self-check)
python mcp_server.py
\end{lstlisting}

Dependencies: Python 3.11+, Lark, Click. Optional: Node.js (for LSP integration), FastMCP (for MCP server mode).

\section{Limitations and Future Work}

\textbf{Limited parser:} Our Lark grammar fails on the full AHK v2 language due to reduce/reduce collisions. The thqby LSP (9,525 lines of TypeScript) is the only complete parser. Integrating it via subprocess or porting its tokenizer state machine to Python would enable more sophisticated analyses.

\textbf{No LLM-driven fixes:} Complex patterns (GUI, labels, DllCall) currently require manual intervention. We plan to use MCP sampling (\texttt{ctx.sample()}) to let the linter request LLM-generated repair suggestions for hard cases.

\textbf{Incremental analysis:} The linter re-scans all files each run. Adding file-hash caching would make it practical as a daemon for IDE integration.

\textbf{Pre-commit hook:} Automating the linter as a pre-commit or CI gate is straightforward with SARIF output and existing CI tooling.

\textbf{Larger corpus:} Testing against a scraped corpus from the AHK forums and GitHub (v1 scripts only) would strengthen the evaluation. We estimate ~10,000 v1 scripts are publicly available.

\section{Conclusion}

We presented \texttt{ahk-lint}, a dedicated migration linter for AutoHotkey v2. Combining rule extraction from community tools with automated analysis, it covers 167 patterns, provides auto-fix for 7/10 primary patterns, and integrates with the MCP protocol for agent self-checking. Our evaluation shows 343 issues across 7 legacy scripts, 32 issues across 5 LLM-generated scripts, and zero false positives on 80+ production v2 scripts. Coverage against the comprehensive community converter stands at 51\%, with a clear path to improvement. AutoHotkey has been neglected by academic tooling research for two decades---this work begins to address that gap.

\subsection*{Data and Code Availability}

\url{https://github.com/sandraschi/ahk-linter}

\subsection*{Acknowledgments}

The authors thank the AutoHotkey community for two decades of contributions that made this work possible. Development was assisted by AI coding tools, as is increasingly common in software engineering research. The rule set was derived from the community-maintained \texttt{AHK-v2-script-converter} by \texttt{mmikeww}.

\begin{thebibliography}{20}
\bibitem{ahk_survey}
\url{https://github.com/search?q=autohotkey+language%3AAHK&type=repositories}. Accessed 2026.
\bibitem{rustfix}
\url{https://github.com/rust-lang/rustfix}. Rust compiler fix suggestions.
\bibitem{stratego}
M. Bravenboer et al., ``Stratego/XT 0.17. A Language and Toolset for Program Transformation,'' \textit{Science of Computer Programming}, 2008.
\bibitem{txl}
J. R. Cordy, ``The TXL Source Transformation Language,'' \textit{Science of Computer Programming}, 2006.
\bibitem{api_migration}
H. Zhong et al., ``Mining API Mapping for Language Migration,'' \textit{ICSE 2010}.
\bibitem{libsync}
S. Negara et al., ``Practical and Effective API Migration via Library Usage Mining,'' \textit{ASE 2020}.
\bibitem{gu2024repair}
J. Gu et al., ``Automated Code Repair with Large Language Models,'' \textit{ICSE 2024}.
\bibitem{mcp_patterns}
Model Context Protocol. \url{https://modelcontextprotocol.io/}. Anthropic, 2024.
\bibitem{dsl_linter}
J. K{\"a}ll{\"a}inen et al., ``Improving DSL Quality with Language-Specific Analysis Tools,'' \textit{Software and Systems Modeling}, 2021.
\bibitem{thqby}
thqby. \texttt{vscode-autohotkey2-lsp}. \url{https://github.com/thqby/vscode-autohotkey2-lsp}, 2023.
\bibitem{converter}
mmikeww. \texttt{AHK-v2-script-converter}. \url{https://github.com/mmikeww/AHK-v2-script-converter}, 2024.
\bibitem{mcp}
Model Context Protocol. \url{https://modelcontextprotocol.io/}, 2024.
\bibitem{lark}
Lark: A modern parsing library for Python. \url{https://github.com/lark-parser/lark}, 2021.
\bibitem{sarif}
SARIF: Static Analysis Results Interchange Format. OASIS, \url{https://sarifweb.azurewebsites.net/}, 2022.
\bibitem{2to3}
Python 2 to 3 converter. \url{https://docs.python.org/3/library/2to3.html}.
\bibitem{gofix}
Go fix tool. \url{https://pkg.go.dev/cmd/fix}.
\bibitem{bibtex}
O. Patashnik, ``BibTeXing,'' 1988.
\bibitem{knuth}
D. E. Knuth, \textit{The Art of Computer Programming}, Addison-Wesley, 1968.
\bibitem{autohotkey}
AutoHotkey v2 Documentation. \url{https://lexikos.github.io/v2/docs/}, 2022.
\bibitem{pyupgrade}
\texttt{pyupgrade}: Upgrade Python syntax. \url{https://github.com/asottile/pyupgrade}.
\end{thebibliography}

\end{document}
